\documentclass[10pt]{article}
\usepackage{setspace}
\onehalfspacing
\usepackage{fontawesome}
\usepackage{svg}
\usepackage[newfloat]{minted}
\usepackage{algpseudocode}
\usepackage{algorithm}
\usepackage[margin=1.2in]{geometry}
\usepackage{amsmath,amsthm,amssymb}
\usepackage{mathtools}
\usepackage[colorlinks=true]{hyperref}
\theoremstyle{definition}
\newtheorem{question}{{\normalfont \faGears}~Question}
\newtheorem{compexercise}{{\normalfont \faLaptop}~Implementation exercise}
\renewcommand{\d}{\mathrm d}


\renewcommand{\mymarks}[1]{\textbf{#1 marks}}
\usepackage{fancyhdr}
\pagestyle{fancy}
\fancyhf{}
\lhead{Name: }
\chead{Final exam}
\rhead{Page \thepage}

\begin{document}

\title{Numerical Analysis: Final exam {\small (\mymarks{50})}}
\author{Urbain Vaes}
\date{Fall $(11111100101)_2$}
\maketitle

\noindent \emph{You are not required to complete every question.
Although the total marks on the exam sum to 55,
your final grade will be calculated out of 50.}
\vspace{.3cm}

\noindent \textbf{Academic integrity pledge:} $\square$ \textcolor{blue}{I certify that I will not give or receive any unauthorized help on this exam,
and that all work will be my own. (Tick $\checkmark$ or copy the sentence on your answer sheet)}.
\vspace{.3cm}

\newpage
\begin{question}
    [Floating point arithmetic, \mymarks{10}]
    True or false?
    \begin{enumerate}
        \item Let $(\placeholder)_2$ denote binary representation.
            Then
            \(
                (1000000000)_2 \times (0.001)_2 = (1000000)_2.
            \)

        \item It holds that
            \(
                (0.111111)_2 + (0.000001)_2 = (1.001)_2.
            \)

        \item
            It holds that $(0.\overline{1100})_2 = \frac{4}{5}$.

        \item
            Let $\varepsilon_{32}$ denote the machine epsilon for the \julia{Float32} format,
            i.e.\ \julia{eps(Float32)}.
            Then the number~3 is representable exactly in this format,
            and the next representable number in the \julia{Float32} format is $3 + 2\varepsilon_{32}$.

        \item
            In Julia, \julia{Float64(3.25) == Float32(3.25)} evaluates to \julia{false}.

        \item
            The spacing (in absolute value) between successive half-precision (\julia{Float16}) floating point numbers is always greater than or equal to \julia{eps(Float16)}.

        \item
            In Julia, define \julia{f(x) = x == Inf ? - x : f(-2x)}.
            Then \julia{f(5.0)} evaluates to \julia{Inf}.

        \item
            There exists $n \in \nat$ such that the machine epsilon for the \julia{Float64} format equals $2^{-n}$.
            In other words, the machine epsilon is an exact negative power of two.

        \item
            In Julia, it holds that \julia{eps()/3 + eps()/3 + 1 > 1} evaluates to \julia{true}.

        \item
            In Julia, the expression \julia{1 + 1024eps() == 2} evaluates to \julia{true}.

        \item
            \textbf{Bonus.}
            Does the Julia expression evaluate to \julia{log2(2 + eps()) == 2} evaluate to \julia{true} or~\julia{false}.

            \textit{Explain briefly:}
            \vspace{1cm}

        \item
            \textbf{Bonus.}
            Does the Julia expression \julia{sqrt(4 + 5eps()) == 2} evaluate to \julia{true} or \julia{false}?

            \textit{Explain briefly:}
            \vspace{1cm}
    \end{enumerate}
\end{question}

\newpage
\begin{question}
    [Interpolation, \mymarks{10}]
    The first part of this question is an open-ended question,
    while the second part is a true/false.

    \begin{enumerate}
        \item (\mymarks{5}, open-ended question)
            Describe the Gregory--Newton method to calculate the interpolation polynomial in the case of equally spaced interpolation nodes at $\{0, 1, 2, \cdots, n\}$ 
            and corresponding data values $\{f(0), f(1), f(2), \cdots, f(n)\}$.
            In particular, address the following points:
            \begin{itemize}
                \item 
                    Write down the expression of the interpolating polynomial;

                \item
                    Discuss the advantages and disadvantages of this interpolation method compared to Lagrange interpolation;

                \item
                    Illustrate the method on a concrete example of your choice.
            \end{itemize}

    \newpage
    \item
        (\mymarks{5}) Are the following assertions true or false?
        \begin{itemize}
            \item 
                Let $f(x) = 2x^4 + x^2 - 3x - 1$,
                and let $f_3 \in \mathcal P_3$ denote interpolation polynomial of $f$ at the four points $x_0 < x_1 < x_2 < x_3$.
                Then
                \[
                    f(x) - f_3(x) = \frac{2}{4!} (x - x_0)(x - x_1)(x - x_2)(x - x_3) \, .
                \]

            \item 
                Let $f \colon \real \to \real$ be the Runge function defined by $f(x) = (1 + 25 x^2)^{-1}$,
                and let $f_n \in \mathcal P_{n}$ denote the interpolation polynomial of $f$ at $n + 1$ equally spaced points in the interval~$[-1, 1]$.
                Then, it holds that
                \[
                    \lim_{n \to \infty} \left( \max_{x \in [-1, 1]} \bigl\lvert f(x) - f_n(x) \bigr\rvert \right) = 0 \, .
                \]

            \item 
                Let $f(x) = \sin(\pi x) + \e^x$,
                and let $f_n$ denote the piecewise linear interpolation of $f$ at $n + 1$ equally spaced points in the interval $[-1, 1]$.
                Then, it holds that
                \[
                    \lim_{n \to \infty} \int_{-1}^{1} \bigl\lvert f(x) - f_n(x) \bigr\rvert^2 \, \d x = 0 \, .
                \]


            \item
                The Runge phenomenon can be mitigated by using Chebyshev nodes instead of equally spaced nodes.

            \item
                In the following code,
                \julia{p} is the interpolating polynomial through the data in \julia{x} and \julia{y}.
                Here, for an array \julia{y} of size $n$, the command \julia{diff(y)} returns the array of size $n-1$
                whose $i$-th entry is equal to \julia{y[i+1] - y[i]}.
        \begin{minted}{julia}
using Plots
x = [0, 1, 2, 3]
y = [1, 2, 1, 2]

function p(x)
    return (y[1]
            + diff(y)[1] * x
            + 1/2 * diff(diff(y))[1] * x^2
            + 1/6 * diff(diff(diff(y)))[1] * x^3
end

plot(p, xlims=(0, 5))
scatter!(x, y)
        \end{minted}
            \item {(\textbf{Bonus, 1 mark})}
                Let $f(n) = 2^n$. Then for all $n \in \nat$, it holds that
                \[
                    f(n) = \sum_{i=0}^{+\infty} \frac{\Delta^i f(0)}{i!} n^{\underline{i}} \, .
                \]
                Justify your answer.
        \end{itemize}
\end{enumerate}
\end{question}

\newpage
\begin{question}
    [Numerical interpolation, \mymarks{10}]
    We wish to derive and implement Boole's rule for numerical integration.
    Boole's rule is a Newton--Cotes quadrature rule based on five equally spaced points in the interval~$[-1, 1]$.
    \begin{equation}
        \label{eq:boole_rule}
        I[u] \coloneq \int_{-1}^{1} u(x) \, \d x
        \approx w_0 u(-1) + w_1 u{\left(- \frac{1}{2}\right)} + w_2 u(0) + w_3 u{\left(\frac{1}{2}\right)} + w_4 u(1)
        \eqcolon \widehat I_5[u].
    \end{equation}

    \begin{enumerate}
        \item
            (\mymarks{5})
            Calculate the weights $w_0, w_1, w_2, w_3, w_4$ to ensure that the quadrature rule $\widehat I_5$ is exact for all polynomials in $\mathcal P_d$,
            for the largest possible value of $d$.
            Include the details of your calculation.

            \textit{\textbf{Hint}: It is useful to require that the quadrature rule is exact for the Lagrange basis polynomials.}

            \textit{\textbf{Hint}: By symmetry, you can expect that $w_0 = w_4$ and $w_1 = w_3$.}

            \textit{Your answer:}
            \vspace{13cm}

        \item
            \mymark{1}
            What is the degree of precision of Boole's rule~\eqref{eq:boole_rule}?

            \textit{Your answer:}

        \newpage
        \item
            (\mymarks{4})
            Write a function \julia{boole(u, a, b)} that returns, using Boole's rule appropriately scaled to the interval $[a, b]$,
            an approximation of the integral
            \begin{align}
                \label{eq:integral_midpoint}
                \int_{a}^{b} u(x) \, \d x \, .
            \end{align}

            \begin{minted}{julia}
    function boole(u, a, b)
        # Write your code here














    end
            \end{minted}
        \item
            (\textbf{Bonus}, \mymarks{1})
            Write a function \julia{composite_boole(u, a, b, N)}
            that returns an approximation of the integral~\eqref{eq:integral_midpoint},
            this time using a composite version of Boole's rule.
            More precisely, the approximation should be obtained
            by partitioning the integration interval $[a, b]$ into $N$ cells,
            and applying Boole's rule within each cell.
            Your function should call the function \julia{boole} you wrote in the previous item.

            \begin{minted}{julia}
    function composite_boole(u, a, b, N)
        # Write your code here











    end
            \end{minted}
    \end{enumerate}
\end{question}


\newpage
\begin{question}
    [Iterative method for linear systems, \textbf{10 marks}]
    Assume that $\mat A \in \real^{n \times n}$ is a nonsingular matrix and that $\vect b \in \real^n$.
    We wish to solve the linear system
    \begin{equation}
        \label{eq:linear_system}
        \mat A \vect x = \vect b
    \end{equation}
    using an iterative method where each iteration is of the form
    \begin{equation}
        \label{eq:iterative_scheme}
        \mat M \vect x_{k+1} = \mat N \vect x_k + \vect b.
    \end{equation}
    Here $\mat A = \mat M - \mat N$ is a splitting of~$\mat A$ such that $\mat M$ is nonsingular,
    and $\vect x_k \in \real^n$ denotes the $k$-th iterate of the numerical scheme.

    \begin{enumerate}
        \item
            (\mymarks{3})
            Let $\vect e_k \coloneq \vect x_k - \vect x_*$,
            where $\vect x_*$ is the exact solution to~\eqref{eq:linear_system}.
            Prove that the error satisfies the recurrence relation
            \[
                \vect e_{k+1} = \mat M^{-1} \mat N \vect e_k.
            \]
            \vspace{5cm}

        \item
            (\mymarks{3})
            Let $L = \norm{\mat M^{-1} \mat N}_{1}$.
            Prove that
            \begin{equation}
                \label{eq:inequality_interpolation}
                \forall k \in \nat, \qquad
                \norm{\vect e_k}_{1} \leq L^k \norm{\vect e_0}_{1}.
            \end{equation}
            \vspace{6cm}

        \item
            \mymark
            In general, is the condition $\norm{\mat M^{-1} \mat N}_{1} < 1$ sufficient
            for convergence when $\vect x_0 \neq \vect x_*$?
            Justify briefly your answer.
            \vspace{3cm}

        \item
            \mymark
            In general, is the condition $\norm{\mat M^{-1} \mat N}_{1} < 1$ necessary
            for convergence when $\vect x_0 \neq \vect x_*$?
            Justify briefly your answer.
            \vspace{3cm}

        \item
            (\textbf{2 marks})
            Let $\vect r_k \coloneq \mat A \vect x_k - \vect b$ denote the residual at iteration~$k$.
            Prove that
            \[
                \vect r_{k+1} = \mat N \mat M^{-1} \vect r_k \, .
            \]
    \end{enumerate}
\end{question}

\newpage
\begin{compexercise}
    [Linear systems, \mymarks{10}]
    Assume that $\mat A \in \real^{n \times n}$ is an \emph{invertible upper triangular} matrix and that $\vect b \in \real^n$.
    We wish to solve the linear system
    \begin{equation}
        \label{eq:linear_system_upper}
        \mat A \vect x = \vect b.
    \end{equation}
    \begin{itemize}
        \item
            (\mymarks{3})
            Solve the linear system~\eqref{eq:linear_system_upper} by backward substitution,
            using pen and paper, in the specific case where~$\mat A$ and~$\vect b$ are given by
            \begin{minted}{julia}
    A = [1.0 2.0 3.0; 0.0 4.0 5.0; 0.0 0.0 6.0]
    b = [6.0; 9.0; 6.0]
            \end{minted}
            \vspace{6cm}
        \item
            (\mymarks{4})
            Suppose that~$\mat A$ is upper triangular.
            Write a function \julia{backward_subs(A, b)} implementing the backward substitution algorithm for solving~\eqref{eq:linear_system_upper} in this case.
            \begin{minted}{julia}
    function backward_subs(A, b)
        n = size(A, 1)
        x = zeros(n)













        return x
    end
            \end{minted}

        \item
            (\mymarks{3})
            How many floating point operations (additions, subtractions, multiplications, divisions) does your implementation of backward substitution perform in order to solve a linear system of size~$n$?
            Include the details of your calculation.

            \vspace{5cm}

        \item
            (\textbf{Bonus, 2 marks})
            Assume now that $\mat A$ is a banded upper triangular matrix with bandwidth~$1$,
            which means that $a_{ij} = 0$ whenever $j > i + 1$.
            Write a function \julia{banded_backward_subs(A, b)} that implements an efficient version of backward substitution for this case.
            \begin{minted}{julia}
    function banded_backward_subs(A, b)
        n = size(A, 1)
        x = zeros(n)
















        return x
    end
            \end{minted}
    \end{itemize}
\end{compexercise}

\newpage
\begin{question}
    [Iterative method for nonlinear equations, \textbf{10 marks}]
    Consider the problem of finding a root of a nonlinear function $f \colon \real \to \real$.
    \begin{itemize}
        \item (\mymarks{2})
            Assume that a numerical scheme of the form
            \[
                x_{k+1} = F(x_k)
            \]
            is employed to find a root of $f$,
            where $F \colon \real \to \real$ is a given function.
            Assume additionally that~$F$ is globally Lipschitz continuous with Lipschitz constant~$L < 1$,
            and that there exists a fixed point~$x_*$ of~$F$.
            First, prove that the fixed point~$x_*$ is unique.
            \vspace{3cm}

        \item
            (\mymarks{3})
            Under the same assumptions as in the previous item,
            prove that the sequence of iterates $(x_k)_{k \in \nat}$ converges to~$x_*$ for any initial guess~$x_0$.
            \vspace{5cm}

        \item
            (\mymarks{3})
            We wish to calculate~$x_* = \sqrt[4]{2}$.
            Propose an iterative scheme~$x_{k+1} = F(x_k)$ that converges to this value.
            Is the convergence of your scheme superlinear? Justify your answer by examining the value of $F'(x_*)$.
    \end{itemize}
\end{question}

\newpage
\begin{compexercise}
    [Iterative methods for nonlinear equations, \mymarks{10}]
    We wish to calculate the point $(x_*, y_*) \in \real^2$ on the parabola $y = x^2$ that is closest to the point $(3, 1)$.
    \begin{itemize}
        \item
            (\mymarks{2})
            Find a function~$f \colon \real \to \real$ such that the root of~$f$ corresponds to the $x$-coordinate of this point:
            \[
                f(x) = \hspace{10cm}
            \]

        \item
            (\mymarks{2})
            Let~$\varepsilon$ denote the dual unit satisfying~$\varepsilon^2 = 0$.
            Calculate $f(x + \varepsilon)$ for $x \in \real$,
            and deduce from the result an expression for $f'(x)$.
            \begin{align*}
                f(x + \varepsilon) &= \hspace{10cm} \\
                \\
                f'(x) &= \hspace{10cm}
            \end{align*}

        \item
            (\mymarks{3})
            Implement the Newton--Raphson method to find an approximation of the root of a general differentiable function~$g \colon \real \to \real$,
            starting from the initial guess $x_0$.
            Use an appropriate stopping criterion, and use automatic differentiation to compute the derivative~$g'$.
            Your function should return only the final approximation of the root.

            \textbf{Hint:}
            Recall that \julia{ForwardDiff.derivative(g, x)} computes the derivative of~$g$ at~$x$.
            \begin{minted}{julia}
    using ForwardDiff
    function newton_raphson(f, x0; tol = 1e-12, maxiter = 10_000)
        # Your code here

















    end
            \end{minted}

        \item
            (\mymarks{2})
            Use your implementation of the Newton--Raphson method (i.e., call the function with appropriate arguments) to find an approximation of the point $(x_*, y_*)$ introduced above,
            and print the result.
            \begin{minted}{julia}
    # Your code here






            \end{minted}

        \item
            (\mymarks{2})
            The secant method is an iterative method for finding a root of a function~$f \colon \real \to \real$,
            where each iteration depends on the two previous iterates,
            according to the formula
            \[
                x_{k+1} = x_k - f(x_k) \frac{x_k - x_{k-1}}{f(x_k) - f(x_{k-1})} \, .
            \]
            Implement the secant method to find an approximation of the root of~$f$,
            starting from the initial guesses $x_0$ and $x_1$.
            \begin{minted}{julia}
    function secant(f, x0, x1; tol = 1e-12, maxiter = 10_000)
        # Your code here
















    end
            \end{minted}
    \end{itemize}
\end{compexercise}


\end{document}
